\section{Understanding Your Project Track}

All three project tracks---Build, Paper, and Competition---are held to equivalent standards of rigor, engineering process, and documentation quality. Your track determines the \emph{form} of your deliverables, not the \emph{level of effort} or the \emph{depth of analysis}. A Paper Track team that delivers a validated simulation is held to the same standard as a Build Track team that delivers a tested prototype. A Competition Track team that places last but executes a strong engineering process will outperform a team that wins the competition but cuts corners on documentation and verification.

The Build Track is well-covered throughout the rest of this guide. This section provides dedicated guidance for Paper Track and Competition Track teams, including sprint-by-sprint adaptations, common pitfalls, and deliverable expectations.


% ═══════════════════════════════════════
\subsection{Paper Track Projects}

\subsubsection{What a Paper Track Project IS and ISN'T}

A Paper Track project produces a complete, verified technical package that a qualified engineer could use to \emph{implement} the solution. The deliverable replaces the physical prototype---it does not replace the engineering rigor.

\textbf{A Paper Track project IS:}
\begin{itemize}[nosep,leftmargin=1.5em]
\item An analytical model, simulation, or computational study
\item A feasibility study with quantitative analysis
\item An optimization or parametric design study
\item A site plan, civil design, or infrastructure analysis
\item A policy analysis grounded in engineering data
\end{itemize}

\textbf{A Paper Track project is NOT:}
\begin{itemize}[nosep,leftmargin=1.5em]
\item A literature review with calculations appended
\item A reduced-effort alternative to building something
\item A report \emph{about} engineering without \emph{doing} engineering
\item A collection of vendor quotes and spec sheets
\end{itemize}

\subsubsection{Sprint-by-Sprint Paper Track Adaptations}

The table below maps key course sprints to their Paper Track equivalents. Sprints not listed follow the standard Build Track guidance.

{\centering\small
\begin{longtable}{c >{\raggedright\arraybackslash}p{5cm} >{\raggedright\arraybackslash}p{5.5cm}}
\toprule
\textbf{Sprint} & \textbf{Build Track Focus} & \textbf{Paper Track Adaptation} \\
\midrule
\endhead
0--1 & Problem definition, stakeholder engagement & Same as Build Track \\
2--3 & Physical test specs, hardware requirements & Requirements include analytical performance targets: convergence criteria, accuracy thresholds, validation benchmarks. Concept generation focuses on modeling approaches, solver selection, boundary conditions. \\
4 (PoC) & Hardware proof-of-concept & Must demonstrate computational goal attainment: converging simulation, validated sub-model, mesh independence study. ``I set up the model'' is not a PoC. \\
5 (PDR) & Preliminary hardware design & Present model architecture, preliminary results, validation plan. Reviewers will ask about assumptions, boundary conditions, mesh sensitivity. \\
6 & Baseline design, long-lead orders & Baseline analytical approach. Begin full-fidelity model development. \\
7--8 (CDR) & Detailed design, fabrication-ready package & Complete model/simulation with preliminary results. Present validation against experimental data or published benchmarks. Must include sensitivity analysis. \\
9--11 & Fabrication, integration, subsystem testing & Full model development, parametric studies, optimization runs. This is your ``fabrication'' phase. \\
12--13 & System-level V\&V testing & Model verification and validation. Independent checks. Documentation of all assumptions. \\
14 (FDR) & Prototype demo, final results & Present complete results, visualizations, recommendations. Deliverable is the technical package, not a prototype demo. \\
\bottomrule
\end{longtable}
}

\subsubsection{Paper Track V\&V}

\begin{faqbox}[What counts as Verification and Validation for a model?]
\textbf{Verification}: Does the model solve the equations correctly?
\begin{itemize}[nosep,leftmargin=1.5em]
\item Mesh convergence or grid independence studies
\item Time-step independence checks
\item Conservation checks (mass, energy, momentum)
\item Unit consistency and dimensional analysis
\item Comparison to known analytical solutions
\end{itemize}

\textbf{Validation}: Do the results agree with reality?
\begin{itemize}[nosep,leftmargin=1.5em]
\item Comparison to experimental data (yours or published)
\item Comparison to published benchmarks
\item Comparison to independent analytical methods
\item Quantified error metrics (RMSE, percent error, confidence intervals)
\end{itemize}

Both are required. A model that runs without errors is \emph{not} verified. A model that matches one data point is \emph{not} validated.
\end{faqbox}

\subsubsection{Paper Track Deliverable Checklist}

Your ``complete technical package'' must contain all of the following:

\begin{enumerate}[nosep,leftmargin=1.5em]
\item \textbf{Model/simulation files} --- documented and version-controlled (Git, not email attachments)
\item \textbf{Input parameter documentation} --- every input justified with a traceable source
\item \textbf{Assumptions log} --- every assumption listed with rationale and sensitivity impact
\item \textbf{Verification evidence} --- convergence studies, conservation checks, analytical comparisons
\item \textbf{Validation evidence} --- comparison to experimental or published data with quantified error
\item \textbf{Results and visualizations} --- publication-quality figures with labeled axes and units
\item \textbf{Recommendations and limitations} --- what the model can and cannot predict, and what should be done next
\item \textbf{User guide} --- step-by-step instructions to reproduce results, including required software, licenses, and computational resources
\end{enumerate}

\begin{warnbox}
Paper Track teams often fall behind because they underestimate the time needed for model validation. Plan to spend at least 30\% of your SDII time on V\&V, not 100\% on running simulations. A beautiful model that cannot be validated is not a successful project.
\end{warnbox}


% ═══════════════════════════════════════
\subsection{Competition Track Projects}

\subsubsection{How Competition Track Differs}

Competition Track teams operate in a unique environment with dual constraints:

\begin{itemize}[nosep,leftmargin=1.5em]
\item Your requirements come from \textbf{two sources}: the course AND the competition rules.
\item You manage \textbf{two timelines}: course sprints AND competition deadlines.
\item Competition placement is \textbf{NOT} part of your course grade---process and deliverable quality are assessed on the same standards as every other team.
\item You still produce \textbf{all course deliverables}: SOW, SDRD, PDRt, CDRt, FDRt.
\end{itemize}

\subsubsection{The Dual Timeline}

You must create a single integrated timeline in Sprint~0 and update it every sprint. When course and competition deadlines conflict, competition deadlines take precedence \emph{operationally}, but course deliverables must still be completed on time.

{\centering\small
\begin{longtable}{>{\raggedright\arraybackslash}p{3.2cm} >{\raggedright\arraybackslash}p{3.5cm} >{\raggedright\arraybackslash}p{4.5cm}}
\toprule
\textbf{Competition Milestone} & \textbf{Typical Timing} & \textbf{Course Gate Alignment} \\
\midrule
\endhead
Registration / team declaration & Early SDI & Sprint 0--1 (SOW) \\
Design report submission & Mid SDI -- Early SDII & Sprint 5--6 (PDR) \\
Technical inspection / safety review & Mid SDII & Sprint 8--10 (CDR through build) \\
Qualifying / practice runs & Late SDII & Sprint 11--12 (V\&V testing) \\
Competition event & End of SDII or later & Sprint 13--14 (FDR) \\
\bottomrule
\end{longtable}
}

\subsubsection{Competition Track Sprint Adaptations}

\begin{description}[leftmargin=1.5em,style=nextline]
\item[Sprint 0--1:] Read the competition rules \emph{cover-to-cover}. Create a compliance matrix mapping every rule to a design requirement. Identify which rules are design-driving vs.\ administrative. Start the SOW with both course and competition scope.

\item[Sprint 2--3:] Requirements capture must include both course requirements AND competition rules. Use MoSCoW prioritization: competition safety rules are \textbf{Must Have}, scoring criteria inform \textbf{Should Have}. Non-scoring rules are still mandatory.

\item[Sprint 4 (PoC):] Demonstrate compliance with the highest-risk competition rule. If the competition has a safety inspection, build toward passing it. Your PoC should retire the biggest technical risk, which is often a competition constraint.

\item[Sprint 5 (PDR):] Present the competition compliance checklist alongside the standard PDR content. Reviewers will ask how you plan to pass tech inspection. Have answers ready.

\item[Sprint 6:] Confirm competition registration. Start travel planning \textbf{NOW} if the competition is in spring---Mines travel approvals require an 8-week lead time. Order all competition-specific components.

\item[Sprint 7--8 (CDR):] Competition compliance checklist should be nearly complete. Travel plan must be finalized. Registration confirmed. Present your integrated timeline showing the path from CDR to competition day.

\item[Sprints 9--11:] Build and test against competition criteria. Prioritize passing tech inspection and safety requirements \emph{before} optimizing for performance. A car that passes inspection but is slow still competes; a fast car that fails inspection does not.

\item[Sprint 12--13:] Final competition preparation. Practice runs under competition conditions. Logistics finalized: vehicle reserved, hotel booked, equipment shipping arranged. Contingency plans documented.

\item[Sprint 14 (FDR):] Present results regardless of competition placement. FDR evaluates your \emph{engineering process and deliverables}, not your ranking. A team that placed 15th but executed excellent engineering judgment, documentation, and V\&V can earn a higher grade than a team that placed 1st with poor process.
\end{description}

\subsubsection{Travel and Logistics}

\begin{tipbox}
Reference the Travel section of this Student Guide for detailed Mines travel policies. Key timeline for competition travel:
\begin{itemize}[nosep,leftmargin=1.5em]
\item \textbf{8 weeks out}: Submit travel request through Mines approval system
\item \textbf{6 weeks out}: Book flights/transportation
\item \textbf{4 weeks out}: Vehicle rental and hotel confirmed
\item \textbf{2 weeks out}: Equipment shipping arranged (if applicable)
\item \textbf{Day of}: Execute the plan
\end{itemize}
International competitions add passport, visa, and customs requirements. See Course Text Ch.~16 for international travel guidance.
\end{tipbox}

\subsubsection{Competition Team Roles}

\begin{protip}
Consider assigning a \textbf{Competition Manager} role alongside your standard Scrum Master, CFO, and other team roles. The Competition Manager is responsible for:
\begin{itemize}[nosep,leftmargin=1.5em]
\item Tracking competition rules, deadlines, and rule changes
\item Managing registration and team enrollment
\item Attending all competition organizer webinars and info sessions
\item Coordinating travel logistics with the CFO
\item Maintaining the competition compliance matrix
\end{itemize}
This role prevents competition logistics from falling through the cracks while the rest of the team focuses on design and build.
\end{protip}

\begin{warnbox}
Do not assume the competition rules are your requirements document. Competition rules tell you what your system must \emph{do}; your SDRD tells you how \emph{well} it must do it and how you will \emph{verify} it. You still need engineering requirements with quantitative acceptance criteria, not just ``pass tech inspection.''
\end{warnbox}


% ═══════════════════════════════════════
\subsection{Hybrid Projects}

Some projects span multiple tracks---for example, a team that builds a prototype \emph{and} enters a competition, or a team that produces both a simulation and a physical test article. If your project is a hybrid, establish the following in Sprint~0:

\begin{enumerate}[nosep,leftmargin=1.5em]
\item \textbf{Clearly define which deliverables serve both tracks and which are unique.} A single SDRD can capture requirements for both streams, but you may need separate V\&V plans.
\item \textbf{Establish a priority framework}: When resources (time, budget, personnel) are constrained, which stream takes precedence? Document this decision and get agreement from your PA and client.
\item \textbf{Watch for scope creep.} Hybrid projects are inherently larger in scope. Resist the temptation to do everything for both streams. Use MoSCoW prioritization aggressively.
\item \textbf{Document the agreement.} Your SOW should explicitly describe the hybrid nature of the project, the priority framework, and what ``done'' looks like for each stream. Get written agreement from your PA and client.
\end{enumerate}

\begin{warnbox}
Hybrid projects are the highest-risk track for scope creep. If you find yourself saying ``we'll just add this for the competition'' or ``we should also simulate this,'' stop and check whether the addition is in your SOW. If not, it requires a formal scope change with PA and client approval.
\end{warnbox}


% ═══════════════════════════════════════
